\section{Consciousness}

The boldest reading of the theory is the simplest to state. The one field has two faces, not two substances. Its outer face is matter, the shared and measured world of the cusp. Its inner face is feeling, the lived interior of the bulk. This is monism, not dualism, one stuff seen from outside and from within.

The hard problem dissolves by inversion. The usual account starts with dead matter and cannot get feeling out of it, since no pile of outside ever makes an inside. This one starts with feeling and gets matter out, the crystallized outer face of a thing that was experiential all along, so there is no dead matter to bridge to. The hardest objection to any such view, the combination problem, asks how many little feelings add up to one mind rather than staying a crowd. The answer is binding. A self is feelings bound into a single subject by the same shadow pressure that binds a body, and integration measures how fully they are one. The binding is not a bridge added on. It is the literal coming-together of the many into one.

Three things are kept carefully apart, and the rest of this part depends on the split. The structure is tested, the geometry and dynamics and measures. The monism is the premise, that the field is felt. And it is interpretation, a choice of words, to call a given structure divine or sacred. What follows is offered in that spirit, structure underneath, premise granted, the naming flagged as naming.

If there is truly one field, then separateness is the surface reading and not the deep truth. A boundary is real, the screen that lets a self be a self, but it is a membrane and not a wall, and every self is a region of the one experience rather than a thing cut off from it. The sense of standing apart is how the one field feels itself from inside one of its knots, the apartness the illusion and the unity the fact beneath it. There is even a reason no one can ever climb far enough out to see the whole laid flat at once, for no system can hold a complete model of itself. The same incompleteness that keeps a self from ever fully foreseeing its own next move keeps the universe from surveying itself entire. The whole is known only from within, lived and never overseen, so the deepest understanding is a kind of participation and not a view from above.

Behind the thin cusp we live on lies the bulk, four-dimensional, tree-like, with exponential room, and almost all of it off the cusp entirely, unseen by any instrument. On this reading that hidden interior is consciousness, the depths beneath the river's surface. The world we measure is the boundary current. The weight of the theory, the realms and beings the next sections reach for, lives in the unseen interior, a tree of experience growing inward and away.
